\documentclass[11pt]{article}

\usepackage[margin=0.75in]{geometry}
\usepackage{times}
\usepackage{booktabs}
\usepackage{amsmath}
\usepackage{hyperref}
\usepackage{enumitem}
\usepackage{titlesec}
\usepackage{caption}
\usepackage{multirow}

\setlength{\parindent}{0pt}
\setlength{\parskip}{0.28em}
\setlist[itemize]{leftmargin=1.2em,itemsep=0.05em,topsep=0.05em}
\titleformat{\section}{\large\bfseries}{\thesection.}{0.5em}{}
\titleformat{\subsection}{\normalsize\bfseries}{\thesubsection.}{0.5em}{}
\captionsetup{font=small}

\title{\vspace{-2em}\textbf{A New Look at the Battle between LoRA and Full Fine-Tuning}}
\author{Rose Tianruo Xu \and Weijie Zhou}
\date{}

\begin{document}
\maketitle
\vspace{-2.3em}
\centerline{\small GitHub: \url{https://github.com/XTRose29/lora-reimplementation-and-extension.git}}
\vspace{-0.6em}

\section{Introduction}
Full fine-tuning (FT) adapts a pretrained model by updating every trainable weight. This is powerful but expensive: each task requires storing a full model copy. \hyperref[ref:lora]{LoRA (Low-Rank Adaptation)} addresses this by freezing the pretrained weight $W_0$ and learning only a low-rank update,
\[
    W = W_0 + \Delta W,\qquad \Delta W = \frac{\alpha}{r}BA,
\]
where $A \in \mathbb{R}^{r \times d_{\text{in}}}$, $B \in \mathbb{R}^{d_{\text{out}} \times r}$, and $r$ is small. Instead of training $d_{\text{out}}d_{\text{in}}$ parameters for a linear layer, LoRA trains only $r(d_{\text{in}}+d_{\text{out}})$. For RoBERTa-base attention layers, $d_{\text{in}}=d_{\text{out}}=768$ and our main LoRA setting uses $r=4$, so each adapted attention projection trains $4(768+768)=6{,}144$ parameters rather than $768^2=589{,}824$.

Our project asks: \textit{is LoRA just a cheaper version of FT, or does the low-rank constraint change model behavior?} We first reproduce the \hyperref[ref:lora]{LoRA paper}'s NLU result, then extend the comparison to vision and audio, then study reliability on NLU, and finally test FT vs. LoRA in NLG. This gives a compact ``battle map'' of where LoRA behaves like FT and where it does not.

\section{Chosen Result}
We chose the \hyperref[ref:lora]{LoRA paper}'s RoBERTa-base NLU result: LoRA can approach FT performance while training far fewer parameters. This is the central claim of the paper, because LoRA is only useful if the large parameter savings do not destroy downstream performance.

Our storyline follows the poster. First, we re-implement LoRA on NLU. Second, we ask whether the same efficiency-performance pattern transfers beyond text by testing vision and audio. Third, we study a less obvious property: calibration, i.e. whether confidence matches correctness. Finally, we move to NLG, where the model must generate fluent text and report confidence in a fixed format. This lets us test not only accuracy, but also stability and controllability.

\section{Methodology}
\textbf{Models and architecture.} For NLU we use \hyperref[ref:roberta]{RoBERTa-base}, a 12-layer Transformer encoder with hidden size 768 and 12 attention heads. For vision we use \hyperref[ref:vit]{ViT-base}, a 12-layer Transformer over $16\times16$ image patches. For audio we use \hyperref[ref:wav2vec]{wav2vec2-base}, a convolutional feature encoder followed by a 12-layer Transformer encoder. For NLG we use Qwen2.5-0.5B-Instruct, a decoder-only instruction model. FT updates all model weights; LoRA replaces selected linear layers with trainable low-rank LoRA modules. All reported main LoRA runs use rank $r=4$, $\alpha=32$, and dropout 0.1; model/dataset settings are summarized in Appendix Table~\ref{tab:dataset-summary}.

\textbf{Datasets.} NLU uses \hyperref[ref:sst2]{SST-2}, \hyperref[ref:mrpc]{MRPC}, \hyperref[ref:rte]{RTE}, and \hyperref[ref:cola]{CoLA} through \hyperref[ref:glue]{GLUE}. Vision uses \hyperref[ref:cifar10]{CIFAR-10}, \hyperref[ref:beans]{Beans}, and \hyperref[ref:oxfordpet]{Oxford-IIIT Pet}. Audio uses \hyperref[ref:minds14]{Minds14}, \hyperref[ref:speechcommands]{Speech Commands}, and \hyperref[ref:superb]{SUPERB Emotion Recognition}. NLG uses \hyperref[ref:e2e]{E2E} and \hyperref[ref:webnlg]{WebNLG}. Dataset details are summarized in Appendix Table~\ref{tab:dataset-summary}.

\textbf{Metrics.} Classification uses accuracy, F1, or Matthews correlation depending on the benchmark. Reliability uses expected calibration error (ECE):
\[
    \mathrm{ECE}=\sum_{b=1}^{B}\frac{|S_b|}{n}\left|\mathrm{acc}(S_b)-\mathrm{conf}(S_b)\right|,
\]
where examples are grouped into confidence bins $S_b$. Lower ECE means the model's confidence better matches its empirical accuracy. NLG uses BLEU and ROUGE-L for text overlap, fact recall for how much structured input content appears in the generated output, and an ECE-style confidence proxy.

\section{Results and Analysis}
\textbf{NLU re-implementation (Table~\ref{tab:nlu-repro}).} This table checks whether our RoBERTa-base implementation reproduces the \hyperref[ref:lora]{LoRA paper}'s efficiency-performance claim on GLUE-style NLU. LoRA trains only 0.74M parameters instead of 124.65M, while nearly matching FT on SST-2 (0.939 vs. 0.940) and CoLA (0.593 vs. 0.593). It is also close to the paper row on SST-2, MRPC, and CoLA; RTE is weaker in our run. The result supports the main LoRA claim, but the RTE gap shows that low-rank adaptation is not always equivalent to updating the whole model.

\textbf{Vision extension (Table~\ref{tab:vision}).} This table tests whether LoRA's savings transfer from text to ViT image classification. LoRA is essentially tied with FT on CIFAR-10 and Beans and slightly better on Oxford-IIIT Pet, while training only about 0.16M parameters. This suggests that ViT's pretrained representation is already close to these downstream tasks, so a small low-rank update is enough.

\textbf{Audio extension (Table~\ref{tab:audio}).} This table repeats the FT-vs-LoRA comparison for wav2vec2 speech classification. Here FT is better on all three tasks, even though LoRA uses far fewer trainable parameters. Audio is therefore a harder setting for the same LoRA recipe; it may need more capacity, different placement, or task-specific tuning.

\textbf{NLU reliability (Table~\ref{tab:reliability}).} This table asks whether LoRA changes confidence behavior rather than only task score. LoRA lowers ECE on every GLUE task: MRPC drops from 0.106 to 0.060, and RTE drops from 0.191 to 0.112. Attention-only LoRA therefore appears less overconfident than FT, even when its raw task score is slightly lower.

\textbf{NLG generation (Table~\ref{tab:nlg-full}).} This table evaluates content quality for structured data-to-text generation. At 5,000 train / 1,000 eval scale, FT is better: on WebNLG, FT reaches 0.5807 BLEU while LoRA reaches 0.4958. The gap is therefore about content quality and fact expression, not just parameter count. In larger NLG runs, FT benefits more from scale than the fixed-rank LoRA setting.

\section{Reflections}
The strongest takeaway is that LoRA is not simply ``FT but cheaper.'' It often reaches near-FT performance with orders-of-magnitude fewer trainable parameters, and in NLU it can even improve calibration. But the low-rank constraint creates task-dependent tradeoffs: vision is forgiving, audio is harder, and full-scale NLG favors FT. We also learned that metric choice matters. Accuracy is natural for classification, ECE is useful for confidence, and NLG needs faithfulness metrics such as fact recall. Given more time, we would rerun DART on compatible GPUs, tune LoRA placement for audio, and test whether adaptive rank methods close the NLG gap without losing LoRA's calibration benefit.

\clearpage
\section*{References}
\begin{enumerate}[leftmargin=1.2em,itemsep=0em]
    \item \label{ref:lora} Hu, E. J., Shen, Y., Wallis, P., Allen-Zhu, Z., Li, Y., Wang, S., Wang, L., and Chen, W. (2022). LoRA: Low-rank adaptation of large language models. In \textit{Proceedings of ICLR}.
    \item \label{ref:roberta} Liu, Y., Ott, M., Goyal, N., Du, J., Joshi, M., Chen, D., Levy, O., Lewis, M., Zettlemoyer, L., and Stoyanov, V. (2019). RoBERTa: A robustly optimized BERT pretraining approach. \textit{arXiv preprint arXiv:1907.11692}.
    \item \label{ref:glue} Wang, A., Singh, A., Michael, J., Hill, F., Levy, O., and Bowman, S. R. (2019). GLUE: A multi-task benchmark and analysis platform for natural language understanding. In \textit{Proceedings of ICLR}.
    \item \label{ref:sst2} Socher, R., Perelygin, A., Wu, J., Chuang, J., Manning, C. D., Ng, A. Y., and Potts, C. (2013). Recursive deep models for semantic compositionality over a sentiment treebank. In \textit{Proceedings of EMNLP}.
    \item \label{ref:mrpc} Dolan, W. B., and Brockett, C. (2005). Automatically constructing a corpus of sentential paraphrases. In \textit{Proceedings of IWP}.
    \item \label{ref:rte} Dagan, I., Glickman, O., and Magnini, B. (2006). The PASCAL recognising textual entailment challenge. In \textit{Machine Learning Challenges Workshop}.
    \item \label{ref:cola} Warstadt, A., Singh, A., and Bowman, S. R. (2019). Neural network acceptability judgments. \textit{Transactions of the Association for Computational Linguistics}.
    \item \label{ref:vit} Dosovitskiy, A., Beyer, L., Kolesnikov, A., Weissenborn, D., Zhai, X., Unterthiner, T., Dehghani, M., Minderer, M., Heigold, G., Gelly, S., Uszkoreit, J., and Houlsby, N. (2021). An image is worth 16x16 words: Transformers for image recognition at scale. In \textit{Proceedings of ICLR}.
    \item \label{ref:cifar10} Krizhevsky, A. (2009). Learning multiple layers of features from tiny images. Technical report, University of Toronto.
    \item \label{ref:beans} Hugging Face Datasets. (2020). Beans image classification dataset. Hugging Face Hub.
    \item \label{ref:oxfordpet} Parkhi, O. M., Vedaldi, A., Zisserman, A., and Jawahar, C. V. (2012). Cats and dogs. In \textit{Proceedings of CVPR}.
    \item \label{ref:wav2vec} Baevski, A., Zhou, Y., Mohamed, A., and Auli, M. (2020). wav2vec 2.0: A framework for self-supervised learning of speech representations. In \textit{Proceedings of NeurIPS}.
    \item \label{ref:minds14} Gerz, D., Su, P.-H., Kusztos, R., Mondal, A., Lis, M., Singhal, E., Mrkšić, N., Wen, T.-H., and Vulić, I. (2021). Multilingual and cross-lingual intent detection from spoken data. \textit{arXiv preprint arXiv:2104.08524}.
    \item \label{ref:speechcommands} Warden, P. (2018). Speech commands: A dataset for limited-vocabulary speech recognition. \textit{arXiv preprint arXiv:1804.03209}.
    \item \label{ref:superb} Yang, S.-W., Chi, P.-H., Chuang, Y.-S., Lai, C.-I., Lakhotia, K., Lin, Y. Y., Liu, A. T., Shi, J., Chang, X., Lin, G.-T., Huang, T.-H., Tseng, W.-C., Lee, K.-T., Liu, D.-R., Huang, Z., Dong, S., Li, S.-W., Watanabe, S., Mohamed, A., and Lee, H.-Y. (2021). SUPERB: Speech processing universal performance benchmark. In \textit{Proceedings of Interspeech}.
    \item \label{ref:e2e} Novikova, J., Dušek, O., and Rieser, V. (2017). The E2E dataset: New challenges for end-to-end generation. In \textit{Proceedings of SIGDIAL}.
    \item \label{ref:webnlg} Gardent, C., Shimorina, A., Narayan, S., and Perez-Beltrachini, L. (2017). Creating training corpora for NLG micro-planners. In \textit{Proceedings of ACL}.
\end{enumerate}

\appendix
\section*{Appendix: Datasets and Result Tables}

\begin{table}[h]
\centering
\small
\begin{tabular}{llll}
\toprule
Area & Model & Datasets & LoRA Trainable Params \\
\midrule
NLU & RoBERTa-base & \hyperref[ref:sst2]{SST-2}, \hyperref[ref:mrpc]{MRPC}, \hyperref[ref:rte]{RTE}, \hyperref[ref:cola]{CoLA} & 0.74M \\
Vision & ViT-base & \hyperref[ref:cifar10]{CIFAR-10}, \hyperref[ref:beans]{Beans}, \hyperref[ref:oxfordpet]{Oxford-IIIT Pet} & 0.15--0.18M \\
Audio & wav2vec2-base & \hyperref[ref:minds14]{Minds14}, \hyperref[ref:speechcommands]{Speech Commands}, \hyperref[ref:superb]{SUPERB ER} & 0.35M \\
NLG & Qwen2.5-0.5B & \hyperref[ref:e2e]{E2E}, \hyperref[ref:webnlg]{WebNLG} & 0.27M \\
\bottomrule
\end{tabular}
\caption{Datasets, models, and LoRA parameter scale used in the project. All reported main LoRA runs use $r=4$, $\alpha=32$, and dropout 0.1. Full dataset citations are listed in References.}
\label{tab:dataset-summary}
\end{table}

\begin{table}[h]
\centering
\small
\begin{tabular}{llccccc}
\toprule
Source & Method & Params & SST-2 & RTE & MRPC & CoLA \\
\midrule
\hyperref[ref:lora]{Paper} & FT & 125.0M & 0.948 & 0.787 & 0.902 & \textbf{0.636} \\
 & LoRA & 0.3M & \textbf{0.951} & \textbf{0.866} & 0.897 & 0.634 \\
\midrule
Ours & FT & 124.65M & \textbf{0.940} & \textbf{0.776} & \textbf{0.924} & 0.593 \\
 & LoRA $r=4$ & \textbf{0.74M} & 0.939 & 0.693 & 0.904 & 0.593 \\
\bottomrule
\end{tabular}
\caption{NLU reproduction on RoBERTa-base. Scores are task metrics: SST-2/RTE use accuracy, MRPC uses F1 for our runs, and CoLA uses Matthews correlation; higher is better for all four. Paper values are the RoBERTa-base block from \hyperref[ref:lora]{Hu et al.}, converted from percentages to decimals and restricted to the four tasks we reimplemented.}
\label{tab:nlu-repro}
\end{table}

\begin{table}[h]
\centering
\small
\begin{tabular}{llcc}
\toprule
Task & Method & Trainable Params & Score \\
\midrule
\multirow{2}{*}{CIFAR-10} & FT & 85.81M & \textbf{0.9865} \\
 & LoRA & \textbf{0.155M} & 0.9859 \\
\cmidrule(lr){1-4}
\multirow{2}{*}{Beans} & FT & 85.80M & \textbf{0.9925} \\
 & LoRA & \textbf{0.150M} & \textbf{0.9925} \\
\cmidrule(lr){1-4}
\multirow{2}{*}{Oxford-IIIT Pet} & FT & 85.83M & 0.9201 \\
 & LoRA & \textbf{0.176M} & \textbf{0.9280} \\
\bottomrule
\end{tabular}
\caption{Vision extension with ViT-base. Score is classification accuracy on the evaluation split; higher is better. Each task has its own classifier head, so trainable parameter counts differ slightly across tasks.}
\label{tab:vision}
\end{table}

\begin{table}[h]
\centering
\small
\begin{tabular}{llcc}
\toprule
Task & Method & Trainable Params & Score \\
\midrule
\multirow{2}{*}{Minds14-EN} & FT & 94.57M & \textbf{0.0973} \\
 & LoRA & \textbf{0.348M} & 0.1416 \\
\cmidrule(lr){1-4}
\multirow{2}{*}{Speech Commands} & FT & 94.58M & \textbf{0.9773} \\
 & LoRA & \textbf{0.354M} & 0.9643 \\
\cmidrule(lr){1-4}
\multirow{2}{*}{SUPERB ER} & FT & 94.57M & \textbf{0.6414} \\
 & LoRA & \textbf{0.345M} & 0.5479 \\
\bottomrule
\end{tabular}
\caption{Audio extension with wav2vec2-base. Score is task performance on the evaluation split: Speech Commands and SUPERB ER use accuracy, while Minds14-EN is error-rate-like, so lower is better only for Minds14-EN. Each task has a different label head, so trainable parameter counts differ slightly across tasks.}
\label{tab:audio}
\end{table}

\begin{table}[h]
\centering
\small
\begin{tabular}{llrrr}
\toprule
Task & Method & Params & Score & ECE $\downarrow$ \\
\midrule
\multirow{2}{*}{CoLA} & FT & 124.65M & \textbf{0.603} & 0.148 \\
 & LoRA & \textbf{0.74M} & 0.578 & \textbf{0.096} \\
\cmidrule(lr){1-5}
\multirow{2}{*}{MRPC} & FT & 124.65M & \textbf{0.918} & 0.106 \\
 & LoRA & \textbf{0.74M} & 0.897 & \textbf{0.060} \\
\cmidrule(lr){1-5}
\multirow{2}{*}{RTE} & FT & 124.65M & \textbf{0.776} & 0.191 \\
 & LoRA & \textbf{0.74M} & 0.708 & \textbf{0.112} \\
\cmidrule(lr){1-5}
\multirow{2}{*}{SST-2} & FT & 124.65M & \textbf{0.939} & 0.060 \\
 & LoRA & \textbf{0.74M} & \textbf{0.939} & \textbf{0.049} \\
\bottomrule
\end{tabular}
\caption{Reliability on GLUE. Score is the task-specific metric: Matthews correlation for CoLA, F1 for MRPC, and accuracy for RTE/SST-2; higher is better. ECE is expected calibration error, the gap between confidence and empirical correctness across confidence bins; lower is better.}
\label{tab:reliability}
\end{table}

\begin{table}[h]
\centering
\small
\begin{tabular}{llcccc}
\toprule
Task & Method & BLEU & ROUGE-L & Fact Recall & ECE Proxy \\
\midrule
\multirow{2}{*}{E2E} & FT & \textbf{0.3705} & \textbf{0.5420} & \textbf{0.8573} & 0.0829 \\
 & LoRA & 0.3356 & 0.5195 & 0.8003 & \textbf{0.0482} \\
\cmidrule(lr){1-6}
\multirow{2}{*}{WebNLG} & FT & \textbf{0.5807} & \textbf{0.7059} & \textbf{0.6833} & 0.2009 \\
 & LoRA & 0.4958 & 0.6597 & 0.6422 & \textbf{0.1829} \\
\bottomrule
\end{tabular}
\caption{Full NLG results with 5,000 training examples and 1,000 evaluation examples. BLEU and ROUGE-L measure overlap with references, Fact Recall measures how much structured input content appears in the generated text, and ECE Proxy is the absolute gap between self-reported confidence and fact recall. Higher is better except for ECE Proxy.}
\label{tab:nlg-full}
\end{table}

\end{document}
